\begin{textsample}{POS Dim 7 – global\_south\_2024\_12 – Score 2.00 – global\_south\_2024\_12/119299247.txt}  \label{ex:f7_pos_152}
Jordan ' s ambivalent stance on Palestine Share Jordan ' s ambivalent stance on Palestine While maintaining a fa? ade of solidarity, Jordan ' s actions often reveal a pragmatic, if controversial, approach to safeguarding its own interests in a volatile region Jordanian reality towards the Palestinians Jordanians have had a wounded, complicated, and manipulative history.
House of Hashim or the Hashemite dynasty rules Jordan with the current King Abdullah II, as the 41st generation of the direct descendants of Prophet Muhammad.
While Hashemites were the Sharif and Emir of Mecca with rule over larger Hejaz, they were violently ousted from the control over the same by the rival Saud family.
This military loss led to the loss of the symbolically vital title of the"Custodian of the Two Holy Mosques"to the Saud family, of modern-day \textbf{Saudi} Arabia.
Hashemites were also to lose monarchial claim of rule on lands of Syria, Lebanon, Palestine and Iraq and were left with only the Jordanian swathes -- which unlike other Arab lands was bereft of oil, natural resources, ports, or any other meaningful lever of power.
They had to be creative in dealing with the situation and remaining relevant and protected.
Geographical contiguity to Palestinian lands led to Jordan demonstrating \textbf{fluctuating} concern that oscillated from supporting the Palestinian cause to willy-nilly, acting against the same.
During the Arab-Israeli War of 1948, Jordan occupied the West Bank and East Jerusalem ( which it annexed in 1950 ), only to lose the same to Israel in the Six-Day War of 1967.
Historians often note the unique culture of"back channels"( facilitating clandestine meetings ) between the Jordanians and Israelis that led to counterintuitive accommodations, even amid war.
In the subsequent Arab-Israel War of 1973 ( Yom Kippur War ), Jordan had curiously avoided joining Arab ranks against Israel.
Jordanian King Hussein had secretly flown to Tel Aviv to warn Israeli Prime Minister Golda Meir of the impending Arab plans to attack Israel.
Ultimately when Jordan was morally forced to support the war against Israel, declassified accounts confirm that it was only to preserve Jordan ' s image in the Arab world and that it had already entered into a tacit agreement with Israel that it would not attack Israeli elements.
It wasn't the first time the Jordanians had worked duplicitously against the Arab or Palestinian cause, as just a couple of years earlier, the infamous and bloody"Black September"( Sep 1970 -- July 1971 ) had occurred.
Jordanian Army had surrounded camps and townships with Palestinians fidayeen of various groups and brutally cut them down to size.
Nearly 3500 were killed in the purge and the PLO ( Palestinian Liberation Organisation ) fighters were expelled to neighbouring Lebanon.
Again, the Jordanians had aligned with Israel before commencing this attack.
Importantly Syrian forces that had intervened on behalf of Palestinian fighters were repulsed leaving 600 Syrians dead and a loss of 120 tanks.
Interestingly, like the historically insincere support of Jordanians towards the Palestinians, one important figure who participated along the Jordanian forces in massacring the Palestinians was a then Pakistani Brigadier ( later dictator ) Zia-ul-Haq, who was on deputation!
In 1987 Jordanian King Hussein and Israeli foreign affairs minister ( later President ) Shimon Peres met secretly to work on a peace treaty which was followed up with Jordanians forsaking their claim on the West Bank, the subsequent year.
By 1994 the two countries signed a formal Peace Treaty ( the second Arab nation after Egypt ) and normalised relations with Israel.
Odd public bickering on issues like Al Aqsa Mosque aside, the top leadership of both countries kept meeting discretely, whilst Jordan kept maintaining the fa? ade of support towards Palestine.
Practically, they engaged and operated as allies.
Post the Hamas terror attack ( 7th Oct 2023 ) on Israel and the subsequent and disproportionate bludgeoning of the Gaza Strip, the Jordanians had no choice but to expel the Israeli ambassador to maintain face in the Arab world.
However, as it was only Iran and its proxies like the Syrian Bashar Al Assad government and the Lebanese Hezbollah who were offering any meaningful resistance to the Israelis -- then the Jordanians did the unthinkable in favour of the Israelis when they shot down Iranian drones attacking Israel.
Practically, the Jordanians weighed in favour of the Israeli military ( all perfunctory concerns towards the beleaguered Palestinians, notwithstanding ).
Jordanians watched silently as Southern Lebanon was attacked and as the Syrian regime was overrun by forces born out of Al Qaeda, ending the last possible faction that had taken on the Israelis in favour of the Palestinians.
For a long, Jordanian King Abdullah had warned against the spectre of a rising"Shia Crescent"( assertion of Shiite rule in Iran, Iraq, Syria and Lebanon ).
He therefore would be presumably pleased to see the decimation of the sectarian rivals, which coincidentally also ended the last modicum of challenge to the Israelis.Beyond vacuous words, the practical conduct of the Jordanians has been suspiciously two-faced on the ostensible redlines as defined by King Abdullah i.e.,"No to abandoning the guardianship of Jerusalem and the holy sites, no to the one-state solution, and no for Jordan to be an alternative homeland for the Palestinians".
The picture has perhaps never been further from the so-called Jordanian redlines, as now, the fate of Palestinians, Jerusalem, and the possibility of a two-state solution, is at its lowest.
History will judge the conduct of Arab monarchies and their lip-serving commitment towards Palestine, but perhaps none as ambidextrously as the Jordanians, who often fancy themselves to be of the highest order. ( The writer, a military veteran, is a former Lt Governor of Andaman \&amp ; Nicobar Islands and Puducherry.
The views expressed are personal )

% matched lemmas: fluctuate, saudi
\end{textsample}
